\section{Actual returns and common estimates}
\label{sec:passage}

We derive the two passages and compose them in trajectory order. For each
fixed finite-dimensional family, all domains and bounds below are chosen
on common compact reductions, with derivative orders prescribed in advance.
State derivatives are taken at fixed physical parameters. The coordinate
$q$ is ordinary and smooth; square roots describe active passages. Rank four
will be needed only for independent physical realization.

\subsection{Actual geometry and a common first return}
\label{subsec:observed-controls}

Let $\lambda$ range over a finite-dimensional smooth parameter family.
The branch fields, switching seams, contacts and transverse sections
are jointly smooth. In each of two disjoint contact boxes,
\begin{equation}\label{eq:matched}
 X^+_\lambda=X^-_\lambda\quad\hbox{on }\{h_\lambda=0\},
 \qquad X^+_\lambda-X^-_\lambda=h_\lambda B_\lambda.
\end{equation}
The switching function has nonzero differential and the inactive branch
is nonvanishing. Its marked orbit has a persistent quadratic contact:
the first seam derivative vanishes and the negative second seam
derivative has a common positive lower bound. The active region lies
above the seam in an inactive flow box.

Transport the inactive first integrals along the regular connections to
one smooth transverse coordinate $q$, with contact levels $\tau_1=0$
and $\tau_2=\delta$. Proposition~\ref{prop:field-recognition} below
justifies this transport for nonidentity connecting maps. In the
$\nu$th box put $z=q-\tau_\nu$. An inactive flow box gives the seam
$z=s^2c(s,\lambda)$ with $c>0$; the signed clock coordinate
$t=s\sqrt{c(s,\lambda)}$ leaves $z$ unchanged and normalizes the seam.
After a positive time change, matching and smooth division give
\begin{equation}\label{eq:passage:local-geometry}
 X^-=\partial_t,\qquad \Sigma=\{z=t^2\},\qquad
 X^+=(1+h\alpha)\partial_t+h\beta\partial_z,
 \qquad h=z-t^2.
\end{equation}
The functions $\alpha,\beta$ are jointly smooth and the active clock is
positive on a common smaller box. This normalization preserves oriented
trajectories and section maps. Throughout, an \emph{actual} return is
the return of those trajectories in the original physical time.

The incoming and outgoing ports are fixed transverse sections identified
by $z$. We require a common neighbourhood of states and parameters in
which positive local inputs make one contact excursion, every entry and
exit lies inside its contact box, and the regular connecting arcs have
no additional seam event. On one fixed incoming interval $I=(-q_0,q_0)$,
including its grazing thresholds, every trajectory follows this itinerary
and has no earlier return to the incoming section. All intermediate port
images lie in common open domains. Section~\ref{sec:sharpness} verifies
these first-return conditions for the local control construction.

\begin{lemma}[A common one-turn return tube]\label{lem:common-tube}
Let a compact embedded nonstationary reference cycle admit fixed tubular
coordinates $(\theta,w)$, with $\theta\in\R/\mathbb Z$, reference $w=0$
and section $\theta=0$. Suppose a finite smooth family of matched fields
has common local Lipschitz bounds on the tube and differs from the
reference by $O(|\lambda|)$. Fix contact cores and transverse ports with
strict margins, and strictly inactive connecting arcs. In every core,
assume positive penetrations have first entry, strict activity and first
moving exit inside the core, with a touch or no event otherwise.
If the angular clock remains positive and these local and port margins
persist, there are one input interval and parameter neighbourhood on which
the chronological composition is the actual first return. It makes one
whole turn, with no extra seam event or earlier section hit.
\end{lemma}

\begin{proof}
Compactness gives $\dot\theta\ge v_*>0$ on a fixed smaller tube for all
small controls. The matched form $X^-_\lambda+(h_\lambda)_+B_\lambda$
is locally Lipschitz with common bounds. Since $w=0$ is the reference
orbit, division by the positive clock gives, up to a possible first exit,
\[
 \left|\frac{dw}{d\theta}\right|\le L|w|+C|\lambda|,
 \qquad |w(\theta)|\le e^L(|w(0)|+C|\lambda|),\quad0\le\theta\le1.
\]
Choose the input interval and parameter neighbourhood once so this strict
bound is inside the tube, core and port margins. A first-exit argument and
continuation give a full turn. The inactive arc margins exclude events
outside the cores; their assumed local entry/exit property covers every
core input, including a tangency. The transverse ports therefore join the
stated maps on common domains. Strict angular advance makes $\theta=1$
the next section hit. No invariance of the smaller input interval by the
return image, or constancy of return time, is needed.
\end{proof}


\subsection{Reversible passages and finite cancellation}
\label{sec:neutral-reference}\label{sec:center-multiplicity}

The two seam intersections of an active leaf form a decreasing
involution. Identifying its entry and exit first will give both the
actual passage and the coefficient relation responsible for cancellation.

\begin{proposition}[The actual quadratic passage]
\label{prop:quadratic-passage}\label{prop:passage:quadratic}
For a local passage with the quadratic geometry~\eqref{eq:passage:local-geometry}
and fixed local transverse ports, there is a jointly smooth
signed companion $g(t,\lambda)$ on a two-sided interval such that
\begin{equation}\label{eq:passage:reversal}
 g(0,\lambda)=0,\qquad g_t(0,\lambda)=1,\qquad
 \iota g\iota=g^{-1},\qquad \iota(t)=-t.
\end{equation}
The physical passage in the ordinary inactive coordinate is
\begin{equation}\label{eq:passage:physical}
 \Gamma(q,\lambda)=q\quad(q\leq0),\qquad
 \Gamma(r^2,\lambda)=g(r,\lambda)^2\quad(r>0).
\end{equation}
Its first entry is $t=-r$ and its first moving exit is $g(r,\lambda)$.
On a common domain,
\begin{equation}\label{eq:passage:quadratic-jet}
 \Gamma(q,\lambda)
 =q+C_1(\lambda)q_+^{3/2}
       +\frac34C_1(\lambda)^2q_+^2
       +q_+^{5/2}L(\sqrt{q_+},\lambda),
\end{equation}
where $L$ is the restriction of a jointly smooth function. The passage
is an orientation-preserving ordinary $C^1$ local diffeomorphism.
For every prescribed finite derivative order, the companion, its
inverse and the smooth root profiles have uniform bounds after a
compact parameter reduction.
\end{proposition}

\begin{proof}
Divide the active field by its positive $t$ component. Let
$K(t,z,\lambda)$ be the $z$ coordinate where its orbit meets $t=0$.
Smooth finite-time ODE dependence gives $K(0,z,\lambda)=z$ and $K_z>0$
on a common smaller box. Since the active orbit equation has zero
$z$ component on the seam, $K_t(t,t^2,\lambda)=0$. Hence the seam
restriction $\widehat K(t,\lambda)=K(t,t^2,\lambda)$ satisfies
\[
 \begin{gathered}
 \widehat K_t=2tE(t,\lambda),\qquad
 E(t,\lambda)=K_z(t,t^2,\lambda)>0,\\
 \widehat K=t^2M(t,\lambda),\qquad
 M(t,\lambda)=2\int_0^1 sE(st,\lambda)\,ds.
 \end{gathered}
\]
Here $E(0,\lambda)=M(0,\lambda)=1$, so
$\mathcal N(t,\lambda)=t\sqrt{M(t,\lambda)}$ is a smooth increasing
local diffeomorphism. The active leaf pairs its seam intersections by
\[
 A=\mathcal N^{-1}\iota\mathcal N,\qquad
 A(0)=0,\quad A'(0)=-1,\quad A^2=\mathrm{id}.
\]

For input $q=r^2>0$, the inactive orbit first reaches the seam at $-r$.
Its active leaf has value $C=\widehat K(-r)>0$ and second intersection
$e=A(-r)>0$. The seam restriction decreases for $t<0$ and increases
for $t>0$; since $K_z>0$, the leaf $K=C$ lies strictly above the seam
precisely between $-r$ and $e$. The positive clock makes them the first
entry and first exit. Their seam derivatives are $2r$ and $-2e$, up
to the positive clock factor. After exit the inactive orbit has $z=e^2$
and lies below the seam; port localization excludes further local events.
Inputs $q<0$ stay below it and $q=0$ only touches it, giving the identity
passage on $q\le0$.

Set $g=A\iota$. Then $A^2=\iota^2=\mathrm{id}$ gives
$\iota g\iota=g^{-1}$ and $g'(0)=1$, while $e=g(r)$ proves
\eqref{eq:passage:physical}. For negative $t$, reversal gives
\begin{equation}\label{eq:passage:negative-companion}
 g(t,\lambda)=-\sqrt{\Gamma^{-1}(t^2,\lambda)}\qquad(t<0).
\end{equation}
This signed continuation records backward active travel; the forward
physical passage on negative $q$ is the identity.

Writing $g(t)=t+\rho t^2+\sigma t^3+O(t^4)$ and comparing cubic
coefficients in $(g\iota)^2=\mathrm{id}$ gives $\sigma=\rho^2$.
This is the classical involution jet relation; see
\cite[Corollary~1 and Proposition~4, arXiv author version]{novaessilva2021}
and \cite[Lemma~16, March 2006 author manuscript]{ofarrell2008}.
Taylor's integral remainder therefore gives
\[
 g(t)=t+\rho t^2+\rho^2t^3+t^4B(t,\lambda)
\]
with $B$ jointly smooth. Squaring at $t=\sqrt q>0$ gives
$C_1=2\rho$ and $C_2=3\rho^2=3C_1^2/4$, proving
\eqref{eq:passage:quadratic-jet}. Its parameter derivatives also give
joint ordinary $C^1$ regularity and $\Gamma_q\to1$ at zero.
Compactness, smooth ODE dependence and a common shrinking give uniform
positivity and the prescribed finite derivative bounds.
\end{proof}

The normalized branch jet determines the leading strength:
\begin{equation}\label{eq:passage:strength-jet}
 C_1(\lambda)=\frac43\beta(0,0,\lambda).
\end{equation}
Indeed, $dz/dt=f(t,z)=h\beta/(1+h\alpha)$ and $K_t=-fK_z$.
On the seam $f=0$, $f_z=\beta$, so $E'(0)=-\beta(0,0)$.
Thus $\mathcal N(t)=t-\beta(0,0)t^2/3+O(t^3)$, and
$\mathcal N^{-1}\iota\mathcal N\iota$ has quadratic coefficient
$2\beta(0,0)/3$; squaring gives \eqref{eq:passage:strength-jet}.
Also, $\partial_q=(2r)^{-1}\partial_r$ bounds the ordinary $j$th state
derivative of $q^{5/2}L(\sqrt q,\lambda)$ by $C_jq^{5/2-j}$ on
$q>0$, for every fixed finite $j$ and after finitely many parameter
derivatives. Its zero extension to $q\le0$ is jointly $C^2$ in state
and parameters; higher ordinary regularity at the threshold is not used.

Write $\Gamma_1,\Gamma_2$ for the two normalized actual passages.
The inactive connections preserve $q$, so chronological continuation gives
\begin{equation}\label{eq:passage:two-contact-map}
 \begin{aligned}
 Q_{m,\delta}(q,p)&=\delta+
 \Gamma_2\bigl(\Gamma_1(q,p,m,\delta)-\delta,p,m,\delta\bigr),\\
 P_{m,\delta}&=R_{p,m,\delta}\circ Q_{m,\delta},
 \end{aligned}
\end{equation}
where $R$ is the actual smooth orientation-preserving regular closing
transition. Write $D=P-q$ and $A_\nu(\lambda)=C_{1,\nu}(\lambda)$.
In an independent observation chart put $p=(a,b,\xi)$ and
$\lambda=(a,b,m,\delta,\xi)$, with $\xi$ in a fixed compact set.
Throughout the main results and actual saturation domains, $m$ denotes the
canonical transported strength sum. Thus
\begin{equation}\label{eq:net-strength}
 A_1(p,m,\delta)+A_2(p,m,\delta)=m,\qquad A(p)=A_1(p,0,0).
\end{equation}
The individual strengths and all higher profiles may depend on every
control. The estimates below only need this identity at $\delta=0$;
alternative aligned extensions are denoted by $m_{\mathrm{alt}}$ below.
We assume $A$ has fixed sign with
$|A(p)|\ge A_*>0$, and after reduction the two passage strengths
remain nonzero and opposite. The observed regular jets are
\begin{equation}\label{eq:regular-jets}
 a=R(0),\qquad b=R'(0)-1,\qquad c=\tfrac12R''(0).
\end{equation}
At the reference $a=b=m=\delta=0$. The aligned curvature
$c_0(p)=\tfrac12R''_{p,0,0}(0)$ has fixed nonzero sign and a uniform
strict margin; after shrinking, $R''$ retains that sign at every relevant
input. Its value may vary with the controls. Upper bounds require no
independence of these observations, as Lemma~\ref{lem:auxiliary-controls}
will show. Physical attainment uses the additional condition
\begin{equation}\label{eq:rank}
 \operatorname{rank}D_\lambda(a,b,m,\delta)=4
 \quad\hbox{at the reference field},
\end{equation}
which makes $(a,b,m,\delta,\xi)$ local parameter coordinates while
retaining all higher coupled profiles.

\begin{lemma}[Finite inverse cancellation]
\label{lem:passage:finite-inverse}
On $\delta=0$, the actual composition has the exact root-profile form
\begin{equation}\label{eq:passage:aligned-jet}
 Q_{m,0}(q)=q+mq_+^{3/2}+\frac34m^2q_+^2
                  +q_+^{5/2}L_0(\sqrt{q_+},a,b,m,\xi).
\end{equation}
Consequently, on the whole slice $m=\delta=0$,
\begin{equation}\label{eq:passage:finite-cancellation}
 Q_{0,0}(q)=q+q_+^{5/2}L_0(\sqrt{q_+},a,b,0,\xi)
\end{equation}
is jointly ordinary $C^2$. Opposite leading strengths suffice; the
complete passage germs need not be inverses.
\end{lemma}

\begin{proof}
On a common smaller domain positive inputs and intermediate images stay
positive. The cubic companion jets and $A_1+A_2=m$ give
\[
 g_2g_1(t)=t+\frac m2t^2+\frac{m^2}{4}t^3
                      +t^4B_0(t,a,b,m,\xi),
\]
with $B_0$ jointly smooth. Squaring proves
\eqref{eq:passage:aligned-jet}; for $q\le0$ both passages are identity.
At $m=0$ the remaining $q_+^{5/2}$ profile has the joint $C^2$
extension proved above.
\end{proof}

The actual activation levels on the incoming section are
\begin{equation}\label{eq:passage:true-thresholds}
 0\quad\hbox{and}\quad a_\delta=\Gamma_1^{-1}(\delta,\lambda),
\end{equation}
because the second local input is $I_2=\Gamma_1(q,\lambda)-\delta$.
The increasing first passage has derivative near one, so its inverse
exists on a common ordinary interval and
\begin{equation}\label{eq:passage:threshold-formula}
 a_\delta=\delta\quad(\delta\le0),\qquad
 a_\delta=\bigl[g_1^{-1}(\sqrt\delta,\lambda)\bigr]^2
                                                \quad(\delta>0).
\end{equation}
Thus $a_\delta$ has the sign of $\delta$, $|a_\delta|\le2|\delta|$,
and for positive separation
$a_\delta=\delta-C_{1,1}(\lambda)\delta^{3/2}+O(\delta^2)$ with a
smooth positive-root profile. It is jointly ordinary $C^1$, and is
generally not ordinary $C^2$ at $\delta=0$ when $C_{1,1}\ne0$.
For $\delta>0$, the intervals $q<0$, $0<q<a_\delta$, $q>a_\delta$
have respectively no excursion, only the first, and both excursions.
For $\delta<0$, the corresponding intervals are $q<\delta$,
$\delta<q<0$, $q>0$, with only the second excursion in the middle.
At either endpoint its excursion becomes a tangency touch. At alignment,
negative inputs are inactive, zero touches both contacts and positive
inputs make both excursions.

On the aligned cancelling slice, the closing expansion gives
\begin{equation}\label{eq:center-displacement}
 D_{0,0}(q,p)=a+bq+c_0(p)q^2
       +q_+^{5/2}M(\sqrt{q_+},p)+q^3S(q,p),
 \qquad c_0(p)=\tfrac12R_{p,0,0}''(0)\ne0.
\end{equation}
The root remainder is jointly $C^2$, with bounds
$|\partial_q^j\partial_p^\beta(q^{5/2}M)|\le C_{j,\beta}q^{5/2-j}$
for $j\le2$. At $a=b=0$ the reference displacement consequently has
$D(0)=D_q(0)=0$ and $D_{qq}(0)=2c_0\ne0$.
The full family is jointly $C^1$. With both strengths nonzero, a fibre
is $C^2$ through both events precisely on $m=\delta=0$: otherwise a
nonzero three-halves coefficient survives at an activation event.
We count event fixed points as distinct points without assigning them
ordinary smooth multiplicities.

A smooth increasing coordinate $x=\phi(q)$, $\phi(0)=0$, $\phi'(0)=s>0$,
changes the reference coefficients to $\widetilde c_0=c_0/s$ and
$\widetilde A_\nu=A_{\nu,0}/\sqrt{s}$. The terms involving $\phi''(0)$
cancel at a fixed point with unit multiplier. The odd root lift
$t\sqrt{\phi(t^2)/t^2}$ preserves reversal; nonlinear coordinate terms
first affect order $5/2$, leaving the forced quadratic coefficient intact.

\subsection{Regular connections and auxiliary controls}

\begin{proposition}[Recognition through regular connections]
\label{prop:field-recognition}\label{prop:passage:field-recognition}
In the common geometry above, let $S_1,S_2$ be smooth increasing regular
transitions, with positive derivative margins and all intermediate images
in their common domains. If the second local level is $\delta_2$, then
\begin{equation}\label{eq:passage:unreduced-return}
 P=S_2\circ\Phi_2\circ S_1\circ\Gamma_1,
 \qquad\Phi_2(z)=\delta_2+\Gamma_{2,\mathrm{loc}}(z-\delta_2)
\end{equation}
has the actual transported form
\begin{equation}\label{eq:passage:recognized-return}
 P(q)=R\{\delta+\Gamma_2(\Gamma_1(q)-\delta)\},
 \qquad\delta=S_1^{-1}(\delta_2),\quad R=S_2\circ S_1.
\end{equation}
The second passage retains reversal and the forced coefficient law.
All four observations are smooth in the physical parameters; rank four
gives a chart $(a,b,m,\delta,\xi)$ with canonical $m=C_{1,1}+C_{1,2}$.
\end{proposition}

\begin{proof}
Set $T(x)=S_1(\delta+x)-\delta_2$ and
$\Gamma_2=T^{-1}\Gamma_{2,\mathrm{loc}}T$.
Then $T(0)=0$, $T'>0$, and
$S_1^{-1}\Phi_2S_1(y)=\delta+\Gamma_2(y-\delta)$, proving the formula.
The signed root lift is
\[
 H(t)=t\sqrt{T(t^2)/t^2},\qquad
 T(t^2)/t^2=\int_0^1T'(st^2)\,ds,\qquad
 g_2=H^{-1}g_{2,\mathrm{loc}}H.
\]
The quotient is smooth, positive and even. Thus $H$ is odd, increasing,
with a smooth bounded inverse on a common root interval; conjugation
preserves reversal. Squaring for $t>0$ gives the actual transported passage.
For $\ell=T'(0)$ its first coefficients are
$C_{1,2}=\sqrt\ell\,C_{1,2,\mathrm{loc}}$ and
$C_{2,2}=\ell C_{2,2,\mathrm{loc}}$, so $C_{2,2}=3C_{1,2}^2/4$.
Nonlinear terms in $T$ first enter the order-$5/2$ remainder.
The transported first integral is $S_1^{-1}$ of the old one; its level,
first entry and moving first exit follow the same trajectories.
The common itinerary therefore gives the chronological first return.
Smooth parameter dependence and the submersion theorem give the observations
and the asserted chart. The reduction itself uses no rank condition.
\end{proof}

The permitted branch and seam perturbations preserve matching, quadratic
contacts, regular arcs and the common itinerary. Finite-time ODE dependence
and the event implicit function theorem give smooth regular transitions
with uniform clock, derivative and domain margins. Opposite strengths and
nonzero closing curvature persist; cancellation is imposed at the reference.
Proposition~\ref{prop:specified-rank} supplies rank four through every such
reference. An alternative extension with $A_1+A_2=m_{\mathrm{alt}}$ on
$\delta=0$ satisfies $m=m_{\mathrm{alt}}+\delta\chi$ by smooth division. Estimates using only
the aligned strength identity also apply to this extension; the saturation
domains always use canonical $m$.

Dependent observations can be treated by independent auxiliary passage
parameters and restriction to the original observations.

\begin{lemma}[Independent auxiliary strength and separation]
\label{lem:auxiliary-controls}
Let a jointly smooth finite-parameter family have the normalized return
\[
 \begin{gathered}
 P_\lambda(q)=R_\lambda\bigl(\widehat\delta(\lambda)+
 \Gamma_{2,\lambda}(\Gamma_{1,\lambda}(q)
                              -\widehat\delta(\lambda))\bigr),\\
 \widehat m(\lambda)=C_{1,1}(\lambda)+C_{1,2}(\lambda),
 \end{gathered}
\]
with the signed companions of Proposition~\ref{prop:quadratic-passage}.
Near a reference with $\widehat m=\widehat\delta=0$, there is a jointly
smooth root-profile family with independent auxiliary parameters
$(m,\delta)$, the same coefficient law and aligned identity
$C_{1,1}+\widetilde C_{1,2}=m$, whose return satisfies
\[
 \widetilde P_{\lambda,\widehat m(\lambda),
                       \widehat\delta(\lambda)}=P_\lambda.
\]
All profiles, inverses and required finite derivatives are bounded on
common smaller domains, and the regular leg remains $R_\lambda$.
\end{lemma}

\begin{proof}
Put $\iota(t)=-t$, $A_2=g_{2,\lambda}\iota$ and
$\mathcal N_2(t,\lambda)=(t-A_2(t,\lambda))/2$.
Since $A_2^2=\id$, $\mathcal N_2A_2=\iota\mathcal N_2$ and
$\mathcal N_2'(0)=1$, so $\mathcal N_2$ is a jointly smooth increasing
local normalizer. For $s=m-\widehat m(\lambda)$ set
\[
 \mathcal N_s(t,\lambda)=\mathcal N_2(t,\lambda)-\frac s4t^2,
 \qquad \widetilde g_2=\mathcal N_s^{-1}\iota\mathcal N_s\iota.
\]
A common shrinking keeps $\mathcal N_s'>0$, and the parameter inverse
function theorem gives its joint smooth inverse. Thus $\widetilde g_2$
is tangent to the identity, satisfies
$\iota\widetilde g_2\iota=\widetilde g_2^{-1}$, and equals $g_{2,\lambda}$
when $s=0$. If $\mathcal N_2=t+d_2t^2+O(t^3)$, the quadratic coefficient
of $\mathcal N_2^{-1}\iota\mathcal N_2\iota$ is $-2d_2$.
The change therefore raises the quadratic root coefficient by $s/2$
and the physical first strength by $s$.
Define $\widetilde\Gamma_2(x)=x$ for $x\le0$ and
$\widetilde\Gamma_2(t^2)=\widetilde g_2(t)^2$ for $t>0$.
Reversibility gives its coefficient law and smooth root remainder, with
$C_{1,1}+\widetilde C_{1,2}=\widehat m+s=m$.
Replace the second passage in the return by $\widetilde\Gamma_2$ and
use the independent translation $\delta$. This recovers the original
return exactly at the observed values. Joint smoothness on compact
reductions gives the uniform derivative and domain bounds. These maps
supply scalar estimates without asserting new independent physical
controls.
\end{proof}

\subsection{Exact divisions and event-uniform values}

The passage profiles also give exact divisions and values across the events,
without restrictions on relative control rates.

\begin{lemma}[Actual division and event-uniform values]
\label{lem:exact-value-division}
On one common root box, for $r>0$, $q=r^2$ and $\delta=r^2v$, the
both-active return has the exact expression
\begin{equation}\label{eq:actual-both-active-division}
 Q-q=mr^3+\tfrac34m^2r^4+r^5L_0+\delta rH,
 \qquad H(0,0,0,p)=\tfrac32A(p),
\end{equation}
where $L_0,H$ are jointly smooth in signed root variables and retain all
actual control dependence. For either sign of $q,\delta$, including the
events,
\begin{equation}\label{eq:quartic-global-value}
 |Q-q-mq_+^{3/2}|
 \le K\{m^2q_+^2+q_+^{5/2}
       +|\delta|\sqrt{q_++|\delta|}+|\delta|q_+^{3/2}\}
\end{equation}
with a common constant. Both assertions apply without observed rank by
Lemma~\ref{lem:auxiliary-controls} and restriction. Physical equality
holds at positive root scale; the signed extension is auxiliary.
\end{lemma}

\begin{proof}
Write $g_\nu(t,\lambda)=tT_\nu(t,\lambda)$ with $T_\nu$ jointly smooth
and $T_\nu(0,\lambda)=1$. For signed $r$ and $\delta=r^2v$, put
\[
 U=T_1(r,\lambda)^2,\quad W=U-v,\quad h=r\sqrt W,
 \qquad K_0=v+WT_2(h,\lambda)^2.
\]
On a common smaller box $W>1/2$. These expressions are smooth and
$Q(r^2)=r^2K_0$ exactly for $r>0$. Since $K_0(0,m,v,p)=1$,
joint Hadamard division gives $K_0-1=r\Phi$. Aligned cancellation gives
$\Phi(r,m,0,p)=m+\tfrac34m^2r+r^2L_0(r,m,p)$; division in the independent
variable $v$ then gives the exact identity
\begin{equation}\label{eq:quartic-Hadamard}
 \Phi=m+\tfrac34m^2r+r^2L_0+vH(r,m,v,p),
 \qquad H(0,0,0,p)=\tfrac32A(p).
\end{equation}
Indeed, the first root jet at $r=0$ is
$A_1+A_2(1-v)^{3/2}$, whose $v$ derivative at cancellation is
$-3A_2/2=3A(p)/2$. The profile arguments contain $\delta=r^2v$,
retaining their coupled dependence. Multiplication by $r^3$ proves
\eqref{eq:actual-both-active-division}.

At fixed $(a,b,m,\xi)$ in a full-rank chart, put
$y=\Gamma_1(q)-\delta$. The total derivative is
\begin{equation}\label{eq:quartic-total-detuning}
 Q_\delta=1-\Gamma_{2,q}(y)
 +\Gamma_{2,q}(y)\Gamma_{1,\delta}(q)+\Gamma_{2,\delta}(y).
\end{equation}
The profile parameter derivatives are $O(x_+^{3/2})$. The closing jets give
\begin{equation}\label{eq:quartic-closing-division}
 R(y)=a+(1+b)y+c(\lambda)y^2+y^3S(y,\lambda).
\end{equation}
Integrate \eqref{eq:quartic-total-detuning} from zero to $\delta$,
holding $m$ and the remaining auxiliary data fixed, or $(a,b,m,\xi)$
in an actual full-rank chart. At an intermediate control $t$,
$y_+=(\Gamma_1(q,t)-t)_+\le Kq_++|t|$. Joint ordinary $C^1$ regularity
and the smooth profiles give, including at active arguments equal to zero,
\[
 |\Gamma_{\nu,q}(x)-1|\le K\sqrt{x_+},\qquad
 |\Gamma_{\nu,t}(x)|\le Kx_+^{3/2}.
\]
Thus $|Q_t|\le K[\sqrt{q_++|t|}+q_+^{3/2}]$, and
\begin{equation}\label{eq:quartic-global-detuning}
 |Q_{m,\delta}-Q_{m,0}|
 \le K|\delta|[\sqrt{q_++|\delta|}+q_+^{3/2}].
\end{equation}
Together with \eqref{eq:passage:aligned-jet} this proves
\eqref{eq:quartic-global-value}, with all coupled profiles retained.
The auxiliary chart gives the same calculation for dependent observations,
followed by restriction to their observed values.
\end{proof}

\begin{corollary}[A field criterion]\label{cor:field-criterion}
Suppose a continuous piecewise smooth planar family has the quadratic
contact geometry and common first-return itinerary of
Proposition~\ref{prop:field-recognition}. If the two transported strengths
are nonzero and cancel at a reference where $R(0)=0$, $R'(0)=1$ and
$R''(0)\ne0$, the common-interval upper bound and the necessary
saturation estimates apply to that family. Independence of the four
observations is required additionally for their independent actual
unfoldings and saturation realizations.
\end{corollary}

\begin{proof}
Combine Proposition~\ref{prop:field-recognition} with
Lemma~\ref{lem:auxiliary-controls} and restrict to the observed controls;
under \eqref{eq:rank}, use the actual observation inverse.
\end{proof}
